\chapter{开放编码形成的范畴、概念及概念内涵}

{
\footnotesize
\begin{longtable}{p{2.5cm}p{3cm}p{9cm}}
\caption{开放编码形成的范畴、概念及概念内涵} \label{tab:open-coding} \\
\toprule
范畴 & 概念 & 概念内涵 \\
\midrule
\endfirsthead
\caption[]{开放编码形成的范畴、概念及概念内涵（续）} \\
\toprule
范畴 & 概念 & 概念内涵 \\
\midrule
\endhead
\midrule
\multicolumn{3}{r}{续下页} \\
\endfoot
\bottomrule
\endlastfoot

具身化文化资本 & 学术技能提升 & 指博士生通过参与海外科研实践与学术训练，在多个技能维度上获得实质性提升，具体包括教学能力的培养、科研方法与技能的习得、研究效能的提升以及语言能力的增强。 \\
 & 学术认知拓展 & 指博士生在认知层面获得的多维度拓展，涵盖跨文化与跨学科认知、研究方法与视角认知、研究方向确立、学术前沿与规则认知、专业基础提升以及学术训练方式认知。 \\
 & 学术思维拓展 & 指博士生在国际学术交流与碰撞中，思维方式发生的转变，包括原有思维视野的突破与拓展，以及对自身思维局限的觉察。 \\
 & 学术规范与经验内化 & 指博士生在留学期间的科研实践与学术互动中，逐渐将学术规范和研究经验转化为自身较为稳定的学术意识与研究习惯的过程，具体包括学术规范的内化、学术经验的积累以及学术品味的建立。 \\
\midrule
客观化文化资本 & 产出型客观化文化资本 & 指博士生主动创造的、以物质或准物质形态存在的学术成果，包括技术平台的搭建、教学课件的制作以及论文、专著、专利等学术产出。 \\
 & 使用型客观化文化资本 & 指博士生在留学期间获取并使用的外部学术资源，包括文献资料、数据与计算资源、实验设备与科研工具等。 \\
\midrule
强关系网络 & 师生关系 & 指博士生与海外导师之间在长期学术合作与日常互动中建立的稳定、紧密的关系。 \\
 & 同门关系 & 指博士生与同一导师团队或同一课题组成员之间形成的强联结关系。 \\
 & 同胞关系 & 指博士生与海外同胞学者或学生之间因共同的语言背景、文化认同与相似处境而建立的紧密的支持性关系。 \\
\midrule
弱关系网络 & 学术会议结识 & 指博士生通过参加学术会议、研讨会等活动建立的弱关系。 \\
 & 社群活动结识 & 指博士生通过线上学术社群、校友群、经人介绍等渠道建立的弱关系。 \\
 & 偶然结识 & 指博士生因空间邻近或偶然机会结识的弱关系。 \\
\midrule
正确科研观构建 & 科研态度 & 指博士生在国际学术流动过程中，对科研目标、研究价值和学术职业形成的较为理性和稳定的理解，表现为更加重视研究质量、学术对话和科研工作的长期价值。 \\
 & 科研专注力 & 指博士生在留学期间，受海外科研环境影响，在科研投入时间与精力上呈现出更高程度的专注与持续性。 \\
 & 内在驱动力 & 指博士生在国际学术流动过程中，基于对科研活动的兴趣与认同而形成的自发性动力，具体体现为科研志趣的确立以及学习动机的增强。 \\
\midrule
自我效能感 & 跨文化沟通信心提升 & 指博士生在异文化学术环境中通过持续交流与实践，逐步克服语言和文化障碍，在跨文化沟通和国际学术交流中建立起较为稳定的自信心。 \\
 & 学术自信心提升 & 指博士生通过海外科研训练、学术互动和阶段性成果积累，对自身学术能力、研究水平和国际学术参与能力形成更为积极的自我评价。 \\
\midrule
乐观 & 积极归因 & 指博士生在面对研究困境与客观限制时，能够将研究受阻的原因归结为外部条件的缺失而非自身能力的不足，避免将客观障碍内化为对自我的否定，维护积极的科研心态。 \\
\midrule
韧性 & 孤独适应能力提升 & 指博士生在独自面对异国陌生环境、脱离原有社会支持网络的过程中，逐步培养出独立生活与独处的心理承受能力。 \\
 & 求助心理障碍克服 & 指博士生在留学过程中逐步消除向导师、同伴等寻求帮助时的心理顾虑，形成更为开放的求助意识与行为习惯。 \\
 & 学术压力下的坚持 & 指博士生在长期高强度的学术压力环境中，能够保持对研究目标的持续投入，展现出在逆境中坚守的意志力与耐受力。 \\
\midrule
留学身份声望 & 留学身份声望 & 指博士生凭借CSC资助留学的经历与身份标签，在国内学术场域中获得的附加声誉与认可。 \\
\midrule
关联性声望 & 导师声誉 & 指博士生通过与知名海外导师的师承关系，间接获取导师学术声誉所带来的声望和知名度。 \\
 & 机构声誉 & 指博士生因在高声誉海外高校或研究机构学习、交流而获得的关联性声望。 \\
\midrule
个人特征 & 目标明确性 & 指博士生在留学前及留学期间对自身学术目标的清晰程度，包括留学前是否设定目标、留学期间是否围绕目标开展活动等。 \\
 & 学科背景 & 指博士生所属学科领域的特点，不同学科在对不同类型学术资源的依赖程度、国际合作惯例、资料获取方式等方面存在差异。 \\
 & 主动性 & 指博士生在留学期间主动寻求学术交流、建立社会关系、参与科研活动的行为倾向。 \\
\midrule
留学环境特征 & 留学国家语言环境 & 指留学目的地国家的主要语言与博士生母语之间的距离，以及该国学术交流所使用的语言生态。 \\
 & 留学机构特征 & 指海外接收机构在学术声誉、科研资源、国际化程度等方面的整体状况。 \\
 & 导师特征 & 指留学期间国外导师的个人特质，如导师指导特征、博士生与导师之间的研究方向相似度及文化差异程度等。 \\
\midrule
流动时长 & 流动时长 & 指博士生在海外学习和研究的持续时间。 \\
\midrule
初职获得 & 留学经历优势 & 指留学经历本身在求职过程中作为优先条件、硬性要求或差异化优势发挥作用，体现为用人单位对留学背景的认可与优先考量。 \\
 & 科研能力证明 & 指通过留学期间的学术产出与科研表现，满足招聘机构的硬性指标和能力期待。 \\
 & 学术网络资源 & 指博士生在留学期间建立的社会关系网络在求职阶段作为个人优势发挥作用。 \\
 & 国际视野素养 & 指留学经历带来的国际化视野、跨文化理解能力和学术格局，成为求职中的软性优势。 \\
 & 职业方向选择 & 指博士生在留学经历的影响下，对自身未来职业路径形成更为清晰的判断。 \\
\midrule
项目获得 & 科研成果支撑 & 指博士生凭借留学期间积累的学术产出，在项目申请中具备更为扎实的前期研究基础，提升申请成功率。 \\
 & 留学身份符号 & 指留学经历本身作为人才项目或基金申请的硬性条件或加分项。 \\
 & 项目内容与方向 & 指留学期间接触的前沿研究激发项目创意点，或留学期间的研究方向延续到回国后的项目申请中。 \\
 & 研究能力证明 & 指留学经历作为科研能力、国际化视野、团队协作经验的证明，写入申请书后增强可信度。 \\
\midrule
职称晋升 & 科研成果优势 & 指留学期间或之后产出的高质量论文为职称晋升提供成果支撑，包括论文数量和质量优势、国际合作发表、高水平期刊发表等。 \\
 & 研究方向延续 & 指留学期间确定的研究方向延续到回国后的科研工作中，形成持续产出，助力职称晋升。 \\
 & 政策制度优势 & 指留学经历在职称评审中享有政策层面的支持，包括减免服务期、替换基层工作经历等。 \\
\midrule
机会拓展 & 访学经历 & 指留学经历带来的后续访学经历。 \\

\end{longtable}
}
